\scanpage{79}
\uline{Corollaire 1} $V$ anneau de val. discrète \inserted{hensélien, corps résiduel $k$}, $K$ corps des fractions, \inserted{$\overline K$ clôture alg.}, $A$ schéma abélien sur $K$, $G$ modèle de Néron de $A$, \inserted{$T$ tore maximal de $G_k$}, \deleted{$\unsure{G}\text{ de }T_\ell(A)$}, $\ell$ nb premier $\ne\operatorname{car}.k$, \deleted{On a $\widetilde{T_\ell(G_k)}|U$} donc
\[
\widetilde{T_\ell(T)}|U\subset\widetilde{T_\ell(G_k)}|U\subset T_\ell(A)
\]
et en prenant les points à valeurs dans $\overline K$
\[
M(\unsure{A})_0^I\subset M(\unsure{A})^I\subset M(A)=T_\ell(A(\overline K))\deleted{\unsure{=T_\ell(A)(\overline K)}}
\]
où $I$ est le groupe d'inertie $\subset\pi=\operatorname{Gal}(\overline K/K)$, et où l'exposant $I$ désigne les invariants sous $I$.
Ceci posé, si $A'$ est une deuxième variété abélienne \inserted{sur $K$}, et $\xi$ une \deleted{\unsure{donnée de}} corr. divisorielle sur $A\times A'$, d'où un accouplement compatible avec $\pi$ donc $I$
\[
\varphi_\xi:M(A)\times M(A')\mathbin{-}T_\ell(\overline K^*)\deleted{\unsure{=T_\ell(\unsure{\mathbb G|\mathbf G}_m)}},
\]
on a : \deleted{\unsure{les relations}} $M(A)^I$ et $M(A')_0^I$ sont orthogonaux, ainsi que $M(A)_0^I$ et $M(A')^I$.

En particulier

\uline{Corollaire 2} Soit $\xi$ une \uline{polarisation} sur $A$, d'où une forme bil. alternée non dégénérée sur
\[
\varphi_\xi:M(A)\times M(A)\mathbin{-}T_\ell(\overline K^*)
\]
\margin{\figure{\begin{tikzpicture}\draw (0,1.5)--(.02,.65);\draw (.02,.2)--(.02,0);\end{tikzpicture}}}
Alors
\[
\varphi_\xi(M^I,M_0^I)=0.
\]
En fait, les deux corollaires sont à priori équivalents, comme on constate aisément $\ldots$

\uline{Remarque} A la limite projective, \inserted{si $k$ est sép. clos}, les homo
\[
{}_{\ell^\nu}A(K)\times{}_{\ell^\nu}T(k)\mathbin{-}{}_{\ell^\nu}k^*\simeq{}_{\ell^\nu}K^*
\]
sont nuls. Cependant les accouplements eux-mêmes ne
\mspage{78}
